% Part: first-order-logic
% Chapter: completeness
% Section: compactness

\documentclass[../../../include/open-logic-section]{subfiles}

\begin{document}

\iftag{FOL}
      {\olfileid[id]{fol}{com}{com}}
      {\olfileid[id]{pl}{com}{com}}

\olsection{Teorema Kekompakan}

Salah satu konsekuensi penting teorema kelengkapan ialah teorema
kekompakan. Teorema kekompakan menyatakan bahwa jika setiap himpunan
bagian \emph{berhingga} dari suatu himpunan !!{sentence} dapat dipenuhi,
maka keseluruhan himpunan tersebut dapat dipenuhi---sekalipun himpunan
itu sendiri tak hingga. Hal ini sama sekali tidak jelas dengan
sendirinya. Sekilas, tampaknya tidak ada yang menutup kemungkinan adanya
himpunan tak hingga dari !!{sentence} yang bersifat kontradiktif, tetapi
yang kontradiksinya baru timbul, boleh dikatakan, dari jumlah yang tak
hingga tersebut. Teorema kekompakan menyatakan bahwa skenario semacam
itu dapat dikesampingkan: tidak ada himpunan tak hingga dari !!{sentence}
yang tidak dapat dipenuhi sementara setiap himpunan bagian berhingganya
dapat dipenuhi. Seperti teorema kelengkapan, teorema ini juga mempunyai
versi yang berkaitan dengan perikutan: jika suatu himpunan tak hingga
dari !!{sentence} mengakibatkan sesuatu, maka suatu himpunan bagian
berhingganya sudah mengakibatkan hal tersebut.

\begin{defn}
  Suatu himpunan $\Gamma$ dari !!{formula} disebut \emph{dapat dipenuhi
  secara berhingga} jika dan hanya jika setiap $\Gamma_0 \subseteq
  \Gamma$ yang berhingga dapat dipenuhi.
\end{defn}

\begin{thm}[Teorema Kekompakan]
\ollabel{thm:compactness}
Pernyataan-pernyataan berikut berlaku bagi setiap himpunan kalimat
$\Gamma$ dan setiap kalimat~$!A$:
\begin{enumerate}
  \item $\Gamma \Entails !A$ jika dan hanya jika terdapat $\Gamma_0
    \subseteq \Gamma$ yang berhingga sedemikian sehingga
    $\Gamma_0 \Entails !A$.
  \item $\Gamma$ dapat dipenuhi jika dan hanya jika dapat dipenuhi
    secara berhingga.
\end{enumerate}
\end{thm}

\begin{proof}
Kita membuktikan (2). Jika $\Gamma$ dapat dipenuhi, maka terdapat
\iftag{FOL}{!!a{structure}~$\Struct{M}$}{!!a{valuation}~$\pAssign{v}$}
sedemikian sehingga $\iftag{FOL}{\Sat{M}}{\pSat{v}}{!A}$ untuk setiap
$!A \in \Gamma$. Tentu saja,
$\iftag{FOL}{\Struct{M}}{\pAssign{v}}$ ini juga memenuhi setiap himpunan
bagian berhingga dari~$\Gamma$, sehingga $\Gamma$ dapat dipenuhi secara
berhingga.

Sekarang andaikan $\Gamma$ dapat dipenuhi secara berhingga. Maka setiap
himpunan bagian berhingga~$\Gamma_0 \subseteq \Gamma$ dapat dipenuhi.
Berdasarkan kesahihan
(\iftag{FOL}{%
  \tagrefs{prfAX/{fol:axd:sou:cor:consistency-soundness},
    prfSC/{fol:seq:sou:cor:consistency-soundness},
    prfND/{fol:ntd:sou:cor:consistency-soundness},
    prfTab/{fol:tab:sou:cor:consistency-soundness}}}{%
  \tagrefs{prfAX/{pl:axd:sou:cor:consistency-soundness},
    prfSC/{pl:seq:sou:cor:consistency-soundness},
    prfND/{pl:ntd:sou:cor:consistency-soundness},
    prfTab/{pl:tab:sou:cor:consistency-soundness}}}),
setiap himpunan bagian berhingga tersebut konsisten. Selanjutnya,
$\Gamma$ sendiri harus konsisten berdasarkan
\iftag{FOL}{%
  \tagrefs{prfAX/{fol:axd:ptn:prop:proves-compact},
    prfSC/{fol:seq:ptn:prop:proves-compact},
    prfND/{fol:ntd:ptn:prop:proves-compact},
    prfTab/{fol:tab:ptn:prop:proves-compact}}}{%
  \tagrefs{prfAX/{pl:axd:ptn:prop:proves-compact},
    prfSC/{pl:seq:ptn:prop:proves-compact},
    prfND/{pl:ntd:ptn:prop:proves-compact},
    prfTab/{pl:tab:ptn:prop:proves-compact}}}.
Berdasarkan kelengkapan (\olref[cth]{thm:completeness}), karena
$\Gamma$ konsisten, himpunan itu dapat dipenuhi.
\end{proof}

\tagprob{FOL}
\begin{prob}
Buktikan (1) dari \olref[fol][com][com]{thm:compactness}.
\end{prob}
\tagendprob

\tagprob{notFOL}
\begin{prob}
Buktikan (1) dari \olref[pl][com][com]{thm:compactness}.
\end{prob}
\tagendprob

\iftag{FOL}{%
\begin{ex}
Dalam setiap model~$\Struct{M}$ dari suatu teori~$\Gamma$, setiap suku
tertutup~$t$ tentu menunjuk !!a{element} dari~$\Domain{M}$. Dapatkah kita
menjamin bahwa sebaliknya juga benar, yakni bahwa setiap !!{element}
dari~$\Domain{M}$ ditunjuk oleh suatu suku tertutup? Dengan kata lain,
adakah teori~$\Gamma$ yang semua modelnya tercakup? Teorema kekompakan
menunjukkan bahwa jawabannya negatif jika $\Gamma$ mempunyai model tak
hingga. Berikut penjelasannya. Misalkan $\Struct{M}$ adalah suatu model
tak hingga dari~$\Gamma$, dan misalkan $c$ adalah !!a{constant} yang
tidak terdapat dalam bahasa~$\Gamma$. Misalkan $\Delta$ adalah himpunan
semua kalimat $\eq/[c][t]$ untuk $t$ suatu suku tertutup dalam bahasa
$\Lang{L}$ dari~$\Gamma$, yakni
  \[
  \Delta = \Setabs{\eq/[c][t]}{t \in \Trm[L],\ t \text{ tertutup}}.
  \]
Suatu himpunan bagian berhingga dari $\Gamma \cup \Delta$ dapat ditulis
sebagai $\Gamma' \cup \Delta'$, dengan $\Gamma' \subseteq \Gamma$ dan
$\Delta' \subseteq \Delta$. Karena $\Delta'$ berhingga, himpunan itu
hanya dapat memuat berhingga banyak suku. Ambil $a \in \Domain{M}$,
!!a{element} dari~$\Domain{M}$ yang tidak ditunjuk oleh satu pun suku
tersebut, dan misalkan $\Struct{M'}$ adalah !!{structure} yang sama
persis dengan $\Struct{M}$
kecuali bahwa $\Assign{c}{M'} = a$. Karena $a \neq \Value{t}{M}$ untuk
setiap~$t$ yang muncul dalam~$\Delta'$, $\Sat{M'}{\Delta'}$. Karena
$\Sat{M}{\Gamma}$, $\Gamma' \subseteq \Gamma$, dan $c$ tidak muncul
dalam~$\Gamma$, juga berlaku $\Sat{M'}{\Gamma'}$. Dengan demikian,
$\Sat{M'}{\Gamma' \cup \Delta'}$ untuk setiap himpunan bagian berhingga
$\Gamma' \cup \Delta'$ dari $\Gamma \cup \Delta$. Jadi, setiap
himpunan bagian berhingga dari $\Gamma \cup \Delta$ dapat dipenuhi.
Berdasarkan kekompakan, $\Gamma \cup \Delta$ sendiri dapat dipenuhi.
Dengan demikian, untuk suatu model berlaku
$\Sat{M}{\Gamma \cup \Delta}$. Setiap $\Struct{M}$ semacam itu merupakan
model dari~$\Gamma$, tetapi tidak tercakup, sebab $\Value{c}{M} \neq
\Value{t}{M}$ untuk setiap suku tertutup~$t$ dari~$\Lang{L}$.
\end{ex}

\begin{ex}
Tinjau !!a{language} $\Lang{L}$ yang memuat !!{predicate}~$<$,
!!{constant} $\Obj{0}$ dan $\Obj{1}$, serta !!{function} $+$, $\times$,
dan $-$. Misalkan $\Gamma$ adalah himpunan semua !!{sentence} dalam
!!{language} ini yang benar dalam !!{structure}~$\Struct{Q}$ dengan
domain~$\Rat$ dan interpretasi yang lazim. Jadi, $\Gamma$ adalah
himpunan semua !!{sentence} dari~$\Lang{L}$ yang benar mengenai bilangan
rasional. Tentu saja, dalam~$\Rat$ (bahkan dalam~$\Real$) tidak terdapat
bilangan~$r$ yang lebih besar daripada~$0$ tetapi lebih kecil daripada
$1/k$ untuk setiap $k \in \PosInt$. Bilangan semacam itu, seandainya
ada, disebut \emph{infinitesimal:} tak nol, tetapi kecil tak hingga.
Teorema kekompakan dapat digunakan untuk menunjukkan bahwa terdapat
model-model dari~$\Gamma$ yang di dalamnya ada infinitesimal. Bahasa kita
tidak mempunyai !!a{function} untuk pembagian (pembagian dengan nol
tidak terdefinisi, sedangkan !!{function} harus diinterpretasikan sebagai
fungsi total). Meskipun demikian, kita tetap dapat menyatakan bahwa
$r < 1/k$, sebab hal ini berlaku jika dan hanya jika $r \cdot k < 1$.
Sekarang misalkan $c$ adalah !!{constant} baru dan misalkan $\Delta$
adalah
\[
\{\Obj{0} < c\}
\cup \Setabs{ c \times \num k < \Obj{1} }{k \in \PosInt}
\]
(dengan $\num{k} = (\Obj{1} + (\Obj{1} + \dots + (\Obj{1} +
\Obj{1})\dots))$ yang memuat $k$ buah~$\Obj{1}$). Untuk setiap himpunan
bagian berhingga~$\Delta_0$ dari~$\Delta$, terdapat suatu~$K$ sedemikian
sehingga setiap !!{sentence} $c \times \num{k} < \Obj{1}$ dalam
$\Delta_0$ mempunyai $k < K$. Jika kita memperluas $\Struct{Q}$ menjadi
$\Struct{Q'}$ dengan $\Assign{c}{Q'} = 1/K$, maka
$\Sat{Q'}{\Gamma_0 \cup \Delta_0}$ untuk sebarang $\Gamma_0 \subseteq
\Gamma$ yang berhingga. Karena itu, $\Gamma \cup \Delta$ dapat dipenuhi
secara berhingga (Latihan: buktikan hal ini secara terperinci).
Berdasarkan kekompakan, $\Gamma \cup \Delta$ dapat dipenuhi. Setiap
model~$\Struct{S}$ dari $\Gamma \cup \Delta$ memuat suatu infinitesimal,
yakni~$\Assign{c}{S}$.
\end{ex}
}{}

\tagprob{FOL}
\begin{prob}
Dalam model standar aritmetika~$\Struct{N}$, tidak ada !!{element}~$k \in
\Domain{N}$ yang memenuhi setiap formula~$\num{n} < x$ (dengan $\num{n}$
ialah $\Obj{0}^{\prime\dots\prime}$ dengan $n$ buah $\prime$). Gunakan
teorema kekompakan untuk menunjukkan bahwa himpunan !!{sentence} dalam
bahasa aritmetika yang benar dalam model standar aritmetika $\Struct{N}$
juga benar dalam !!a{structure}~$\Struct{N'}$ yang memuat !!a{element}
yang \emph{memang} memenuhi setiap formula $\num{n} < x$.
\end{prob}
\tagendprob

\iftag{FOL}{
\begin{ex}
Kita mengetahui bahwa logika orde pertama dengan !!{identity} dapat
menyatakan bahwa ukuran !!{domain} harus sekurang-kurangnya sebesar suatu
ukuran minimum: kalimat~$!A_{\ge n}$ (yang mengatakan ``terdapat
sekurang-kurangnya $n$ objek yang berbeda'') hanya benar dalam struktur
yang $\Domain{M}$-nya mempunyai sekurang-kurangnya~$n$ objek. Jadi, jika
kita mengambil
  \[
  \Delta = \Setabs{!A_{\ge n}}{n \ge 1}
  \]
maka setiap model dari $\Delta$ harus tak hingga. Dengan demikian, kita
dapat menjamin bahwa suatu teori hanya mempunyai model tak hingga dengan
menambahkan~$\Delta$ kepadanya: model-model dari $\Gamma \cup \Delta$
ialah tepat semua model tak hingga dari~$\Gamma$.

Jadi, logika orde pertama dapat menyatakan ketakhinggaan. Akan tetapi,
teorema kekompakan menunjukkan bahwa logika ini tidak dapat menyatakan
keberhinggaan. Andaikan, sebaliknya, suatu himpunan kalimat $\Lambda$
dipenuhi dalam tepat semua !!{structure} berhingga. Maka $\Delta \cup
\Lambda$ dapat dipenuhi secara berhingga. Mengapa? Andaikan $\Delta'
\cup \Lambda' \subseteq \Delta \cup \Lambda$ berhingga dengan $\Delta'
\subseteq \Delta$ dan $\Lambda' \subseteq \Lambda$. Definisikan $n$
sebagai 1 jika komponen pertama itu kosong, dan selain itu sebagai bilangan
terbesar sedemikian sehingga $!A_{\ge n} \in \Delta'$. Karena dipenuhi dalam
setiap struktur berhingga, $\Lambda$ mempunyai model~$\Struct{M}$ yang
memiliki berhingga banyak !!{element}, tetapi jumlahnya $\ge n$. Dengan
demikian, $\Sat{M}{\Delta' \cup \Lambda'}$. Berdasarkan
kekompakan, $\Delta \cup \Lambda$ mempunyai model tak hingga, yang
bertentangan dengan asumsi bahwa $\Lambda$ hanya dipenuhi dalam
!!{structure} berhingga.
\end{ex}
}{}

\end{document}
